\section{Strict vanishing from one local decomposition group}
\label{sec:one-prime-proell-original-strict}

\paragraph{The original scalar action is retained.}
This argument uses the maximal pro-$p$ quotient of the actual global
Galois group.  It applies when the coefficient action is pro-$p$ and
unramified away the unique place above $p$.
It includes wild two-power characters and arbitrary ramification of
the original coefficient ring.  The finite torsion module containing
a cocycle is enlarged to contain its local coboundary witness as well.
This supplies the necessary comparison without an invariant-divisibility
assumption or an integral character projector.

For a totally imaginary number field $L$, a finite place $v$, and a
discrete module $D$ unramified outside $v$, write
\[
 \begin{split}
 \mathcal S_v(L,D)=\ker\bigl(&H^1_{\rm cts}(G_L,D)
       \longrightarrow H^1(G_{L_v},D)\\
       &\times\prod_{w\ne v,\ w\text{ finite}}H^1(I_w,D)\bigr).
 \end{split}\tag{G182.1}
\]
An embedding of algebraic closures fixes the local map; changing it
conjugates its image and does not change any conclusion below.

\begin{lemma}[One-prime local surjectivity and cohomology]
\label{G182-one-prime-lemma}
Let $p$ be a rational prime.  Suppose that $L$ is totally imaginary,
has exactly one place $v$ above $p$, and
$\operatorname{Cl}(L)[p]=0$.
Let $\Gamma=G_{L,\{v\}}$ be the Galois group of the maximal extension
unramified outside $v$, and let $P=\Gamma^{(p)}$ be its maximal
pro-$p$ quotient.  Then
\[
                  G_{L_v}\longrightarrow P
                       \quad\text{is surjective}.
 \tag{G182.2}
\]
Let $E/\mathbb Q_p$ be finite, with its original ring of integers $O$
and maximal ideal $\mathfrak m$.  For every continuous character
$a:\Gamma\to1+\mathfrak m$, put $D=(E/O)(a)$.
Then restriction is injective:
\[
 H^1_{\rm cts}(\Gamma,D)\lhook\joinrel\longrightarrow
                      H^1(G_{L_v},D).
 \tag{G182.3}
\]
Consequently $\mathcal S_v(L,D)=0$ when $a$ is viewed as an
unramified-outside-$v$ character of $G_L$.
\end{lemma}

\begin{proof}
\proofstep{1. The finite $p$-group fact.}
Every maximal proper subgroup $M$ of a finite nontrivial $p$-group $Q$
is normal and has index $p$.  Here is an induction proof.
The class equation shows that the center $Z$ of $Q$ is nontrivial.
If $Z\not\subset M$, maximality gives $Q=MZ$.
The centrality of $Z$ makes $M$ normal, and the quotient $Q/M$ has
no proper nontrivial subgroup; as a $p$-group it therefore has order $p$.
If $Z\subset M$, the subgroup $M/Z$ is maximal proper in $Q/Z$,
so induction gives the same conclusion for $M$.
The cyclic group of order $p$ starts the induction.
Every proper subgroup of a finite group is contained in a maximal one,
by taking a subgroup of largest order among the proper overgroups.

\proofstep{2. A proper closed local image would give an unramified character.}
The image $H$ of $G_{L_v}$ in $P$ is compact, hence closed.
If $H\ne P$, choose $g\notin H$.
An open normal subgroup $U$ of $P$ can be chosen so that
$gU\cap H=\varnothing$: open normal subgroups give a neighborhood
basis in a profinite group and $P\setminus H$ is open.
Thus the image of $H$ in the finite $p$-group $Q=P/U$ is proper.
Step 1 produces a nonzero character
\[
                 \Gamma\longrightarrow Q\longrightarrow C_p
 \tag{G182.4}
\]
which vanishes on the local image.
Its degree-$p$ cyclic extension is unramified away $v$, by the
definition of $\Gamma$, and split at $v$, because the whole
decomposition group maps to zero.  It is therefore everywhere
unramified.  There are no real places in $L$.

The Hilbert class-field reciprocity isomorphism identifies every
everywhere unramified abelian extension with a quotient of
$\operatorname{Cl}(L)$.
The hypothesis $\operatorname{Cl}(L)[p]=0$ forbids the quotient $C_p$.
This contradiction proves (G182.2).
The precise retained class-field input is Sharifi,
\emph{Iwasawa Theory}, Theorem 1.1.11 and its preceding definition
of the Hilbert class field.  This is the standard reciprocity theorem
already used in the earlier class-field arguments; the present proof
does not claim a new proof of global class field theory.

\proofstep{3. Finite semidirect images retain the local witness.}
Choose a uniformizer $\pi$ of the original $O$.
For each $m\geq1$, the additive module
\[
       D_m=D[\pi^m]\simeq O/\pi^m O
 \tag{G182.5}
\]
is a finite $p$-group.  The displayed isomorphism is of $O$-modules,
with the scalar action transported to the right-hand side.
The action of $1+\mathfrak m$ on $D_m$
has finite $p$-group image $U_m$.
Indeed the filtration by $1+\pi^iO$, $i\geq1$, has successive
quotients isomorphic to additive subgroups of the residue field,
and $1+\pi^mO$ acts trivially on $D_m$.
Thus $D_m\rtimes U_m$ is a finite $p$-group.

Let $z:\Gamma\to D$ be a continuous cocycle whose local cohomology
class is zero.  Compactness of $\Gamma$ and discreteness of $D$
make the image of $z$ finite.
Choose $d\in D$ such that
\[
       z(g)=(a(g)-1)d\qquad(g\in G_{L_v}).
 \tag{G182.6}
\]
Now choose $m$ large enough that \emph{both} the image of $z$ and
the element $d$ lie in $D_m$.
The cocycle identity says precisely that
\[
 \rho_z:\Gamma\longrightarrow D_m\rtimes U_m,
              \qquad g\longmapsto(z(g),a(g)|_{D_m})
 \tag{G182.7}
\]
is a continuous homomorphism, where
$(x,u)(y,w)=(x+uy,uw)$.
Its finite image is a $p$-group, so it factors through $P$ by the
definition of the maximal pro-$p$ quotient.
Equation (G182.2) implies that its local image equals its global image.
For every $g\in\Gamma$, choose a local element with the same
$\rho_z$-image.  Equation (G182.6), evaluated on that element,
then gives $z(g)=(a(g)-1)d$ globally, because $d\in D_m$.
Thus $z$ is a global coboundary, proving (G182.3).

\proofstep{4. Exact passage from the full strict group.}
Suppose now that $z:G_L\to D$ is strict in (G182.1).
At every $w\ne v$, the coefficient is unramified, so inertia acts
trivially on $D$.  A zero restriction in $H^1(I_w,D)$ therefore
means $z(i)=0$ for every $i\in I_w$: every inertia coboundary is zero.
Choose a finite module $D_m$ containing the cocycle image and its
local witness, as in Step 3.
The corresponding homomorphism (G182.7), now defined on $G_L$,
is trivial on every such inertia group.  Its closed normal kernel
therefore contains their closed normal closure.
It factors through the actual quotient $\Gamma$.
The local witness is unchanged, so (G182.3) proves the strict class
zero.  This argument uses the full local condition at $v$.
\end{proof}

\begin{theorem}[AC-G182: original strict zeros at the one-prime coefficients]
\label{AC-G182}
Put $K=\mathbb Q(\omega)$.
For each rational prime $p=3$ or $p\equiv2\pmod3$, the field $K$ has one place above
$p$ and satisfies the preceding lemma.  Thus every continuous scalar
character $a:G_K\to1+\mathfrak m_O$ unramified away $p$ gives
\[
                       \mathcal S_p(K,(E/O)(a))=0
 \tag{G182.8}
\]
on its original finite coefficient extension $E/\mathbb Q_p$.
Here $\mathcal S_p$ means (G182.1) at the unique $p$-place.

In particular the following two actual families have strict zero.

\emph{(i) Coefficient two, every pure two-power label.}
Let $\delta$ be any finite Dirichlet character of two-power conductor,
including the trivial character.  Suppose the original $O_2$ contains
its values.  For every $e\geq1$ and every continuous
$\kappa:G_K\to1+\mathfrak m_{O_2}$ unramified away two,
\[
 \mathcal S_2\bigl(K,(E_2/O_2)
                    (\delta\kappa^{-1}c_2^{-e})\bigr)=0.
 \tag{G182.9}
\]
There is no conductor or parity exception.
For an actual page $h\geq1$, all $2^{h+1}$ characters modulo
$2^{h+2}$ supply distinct original-coefficient labels at every weight.

\emph{(ii) Trivial finite labels at inert odd coefficient primes.}
Let $\ell\geq5$, $\ell\equiv2\pmod3$, and let $E_\ell/\mathbb Q_\ell$
be any original finite coefficient extension.
For every positive $e$ divisible by $\ell-1$, and every continuous
$\kappa:G_K\to1+\mathfrak m_{O_\ell}$ unramified away $\ell$,
\[
 \mathcal S_\ell\bigl(K,(E_\ell/O_\ell)
                         (\kappa^{-1}c_\ell^{-e})\bigr)=0.
 \tag{G182.10}
\]
At coefficient three the same assertion holds for every positive $e$,
because $c_3|_{G_K}$ already has image in $1+3\mathbb Z_3$.
This last trivial-label row is retained in the existing window unions.
\end{theorem}

\begin{proof}
\proofstep{5. The base field has no class obstruction.}
The quadratic integral-basis calculation gives
$\mathcal O_K=\mathbb Z[\omega]$.
For a complex number written $x+y\omega$ with real $x,y$, choose
integers $a,b$ with $|x-a|,|y-b|\leq1/2$.
Writing $u=x-a$ and $w=y-b$ gives
\[
        |u+w\omega|^2=u^2-uw+w^2\leq\tfrac34<1.
 \tag{G182.11}
\]
Applying this to a quotient of two nonzero Eisenstein integers proves
Euclidean division for the absolute norm.
Hence $\mathcal O_K$ is a principal ideal domain and
$\operatorname{Cl}(K)=0$.
For $p\ne3$, the polynomial $X^2+X+1$ has a root modulo $p$
exactly when the multiplicative group has a nontrivial cube root of
unity, equivalently $p\equiv1\pmod3$.
Thus $p\equiv2\pmod3$ is inert and gives exactly one place.
The discriminant is $-3$, and three is ramified with a unique place.
The lemma proves (G182.8) for all the stated primes.

\proofstep{6. All actions in the two actual families are pro-prime.}
At coefficient two, the reduction of every two-power root of unity
is one, since a finite residue field has multiplicative group of odd
order.  Every value of the two-adic cyclotomic character is a rational
two-adic unit and also reduces to one.
The factor $\kappa$ reduces to one by hypothesis.
Thus the scalar action in (G182.9) lies in $1+\mathfrak m_{O_2}$;
all its factors are unramified away two.
This remains true when $O_2$ contains additional odd-order roots of
unity: its coefficient ring has trivial Galois action and only the
displayed scalar character acts on the module.
The lemma applies directly, with the original $O_2$.

For odd $\ell$, Fermat's theorem gives
$c_\ell(g)^{\ell-1}\equiv1\pmod\ell$.
If $\ell-1\mid e$, the scalar in (G182.10) therefore lies in
$1+\mathfrak m_{O_\ell}$ and is unramified away $\ell$.
The one-prime lemma applies for inert $\ell$.
At three the inclusion $\mu_3\subset K$ makes the cyclotomic
reduction trivial on $G_K$, proving the stated all-weight variant.
No local invariant vanishing or divisibility is needed for any of
these applications: the local coboundary witness is carried through
the same finite semidirect quotient.

\proofstep{7. The original host, coefficient module and full local conditions.}
The actual host (FG.3) contains $\mathbb Q(\zeta_{2^{h+2}})$.
Its pure two-power characters remain distinct on $G_K$, since a
rational character trivial on $G_K$ is trivial or $\chi_{-3}$,
and the latter has conductor three.
For $h\geq1$ their orders divide $2^h$, hence the actual exponent
$N=3\cdot2^h$ in (FC.1).
Their values already lie in the original completion $O_2$.
The trivial finite character at any other coefficient prime also
uses its original $O_\ell$.
The canonical scalar identity is
\[
 (O_\ell\otimes_{\mathbb Z_\ell}\mathbb Q_\ell/\mathbb Z_\ell)
                    (\eta)(-e)
       =(E_\ell/O_\ell)(\eta c_\ell^{-e}).
 \tag{G182.12}
\]
Both sides use the same finite free $\mathbb Z_\ell$-module and the
same Galois action; no scalar extension is made.
Set $\kappa=1$ for the actual labels.

Inflating a cocycle from the original $G_{K,\Sigma_\ell}$ to $G_K$
preserves its full local zero at every place above $\ell$ and its
inside-set inertia conditions.  Outside that set the inflated cocycle
already kills inertia, because both the coefficient and the quotient
are unramified there.  Thus its class lies in (G182.1) and is zero.
If an inflated cocycle is a global coboundary, its original cocycle
is the same coboundary on the quotient, so inflation is injective.
This verifies the literal (FC.2) strict groups, including every added
ramified prime of the growing host, on their original rings.
\end{proof}

\paragraph{Exact sizes available for the subsequent window union.}
For integers $T\geq1$ and $h\geq1$, coefficient two supplies exactly
\[
                         2^{h+1}(T-1)
 \tag{G182.13}
\]
actual pure-character/positive-weight pairs.
At inert primes $\ell\geq5$ distinct from $269$, the newly specified
trivial-character pairs have total count
\[
 J(T)=\sum_{\substack{5\leq\ell\leq T\\
                     \ell\text{ prime},\ \ell\equiv2\bmod3\\
                     \ell\ne269}}
                     \left\lfloor\frac{T-1}{\ell-1}\right\rfloor.
 \tag{G182.14}
\]
Indeed the positive multiples of $\ell-1$ below $T$ have that count;
one can occur only when $\ell\leq T$.  No prime-counting estimate
is used.  Coefficient two is distinct from every prior omitted
coefficient prime.  At primes at least five other than $269$, the
old G175 rows have a nontrivial high two-power factor, while the
new rows have trivial finite character; hence those sets are disjoint.
The coefficient-$269$ trivial row and the coefficient-three trivial
row are already included in the preceding unions and are not counted
again in (G182.14).

\paragraph{Source and geometric scope.}
The class-field input is the retained Hilbert class-field reciprocity
theorem, freshly checked at Sharifi Theorem 1.1.11:
\url{https://www.math.ucla.edu/~sharifi/notes/iwasawa-ch01.html}.
The maximal extension unramified outside a specified set and its
Galois group use Lemma 1.4.2, with its complete printed proof, and
Definition 1.4.3 of the same source.
All finite $p$-group, closed-image, semidirect-cocycle and original-ring
arguments needed here are proved above.
The finite-cover interpretation is the retained SGA~1, Expos\'e~V,
sections 4 and 7 dictionary: the finite semidirect images give finite
Galois covers, and local restriction is their decomposition action.
The Tate coefficient is still the one in (FC.2).
This is no assertion about a nontrivial cubic or conductor-seven
factor, an auxiliary ramified quadratic twist, or a split coefficient
prime.  The numerical positive-index cost and the same-window
comparison still require the single full-union (FL.4) estimate and
(EU.2) in their original normalizations.
No all-prime numerical remainder, mixed realization or H04.4
conclusion is supplied here, and no Lean execution is claimed.
